\chapter*{Приложение А. Параметры экспериментов и структура прототипа}
\addcontentsline{toc}{chapter}{Приложение А. Параметры экспериментов и структура прототипа}
\label{app:params}

\section*{А.1. Гиперпараметры residual RL}
\addcontentsline{toc}{section}{А.1. Гиперпараметры residual RL}

    В таблице~\ref{tab:hyperparams} приведены гиперпараметры SAC и PPO для RL-дообучения VLA на LIBERO.

\begin{table}[htbp]
\caption{Гиперпараметры RL-дообучения VLA (Stable-Baselines3).}
\label{tab:hyperparams}
\centering
\footnotesize
\setlength{\tabcolsep}{5pt}
\renewcommand{\arraystretch}{1.15}
\begin{tabularx}{\textwidth}{@{}>{\raggedright\arraybackslash}Xcc@{}}
\toprule
\textbf{Параметр} & \textbf{SAC (residual)} & \textbf{PPO (STARE)} \\
\midrule
Скорость обучения $\alpha$       & $3 \times 10^{-4}$ & $3 \times 10^{-4}$ \\
Размер скрытых слоёв residual   & 256 / 256          & 256 / 256 \\
Функция активации                & ReLU               & tanh \\
Размер буфера (SAC)              & $10^6$             & -- \\
Размер пакета                    & 256                & 2048 \\
$\tau$ (мягкое обновление SAC)   & 0,005              & -- \\
$n_{\text{epochs}}$ (PPO)        & --                & 10 \\
Горизонт эпизода $H$             & \multicolumn{2}{c}{220 шагов (LIBERO)} \\
Бюджет обучения                  & \multicolumn{2}{c}{300\,000 шагов/набор} \\
Число seed-ов                    & \multicolumn{2}{c}{3 (7, 13, 42)} \\
Заморозка VLA                    & полная (SAC)       & последние 2 слоя (PPO) \\
\bottomrule
\end{tabularx}
\end{table}

\section*{А.2. Стадийная функция вознаграждения}
\addcontentsline{toc}{section}{А.2. Стадийная функция вознаграждения}

    Четыре фазы манипуляции с равными весами $w_k = 0{,}25$:
\begin{enumerate}
    \item \textbf{Reach} -- расстояние захвата до объекта $< 0{,}05$~м;
    \item \textbf{Grasp} -- контакт с объектом и закрытие захвата;
    \item \textbf{Transport} -- объект перемещён, расстояние до цели $< 0{,}15$~м;
    \item \textbf{Place} -- успех эпизода $S(T_i) = 1$, бонус $\beta = 1{,}0$.
\end{enumerate}

\section*{А.3. Протокол оценки VLA на LIBERO}
\addcontentsline{toc}{section}{А.3. Протокол оценки VLA на LIBERO}

    10 эпизодов/задачу (100 эпизодов/набор); 3 seed среды (7, 13, 42); разрешение 224$\times$224; greedy-инференс. Чекпоинты: OpenVLA-OFT (\texttt{moojink/openvla-7b-oft-finetuned-libero-*}), UniVLA, $\pi_0$ и $\pi_0$-fast (openpi), CogACT, SmolVLA, SpatialVLA, Octo-Small, RT-1-X (симуляционный адаптер).

\section*{А.4. Сводная таблица экспериментальных серий}
\addcontentsline{toc}{section}{А.4. Сводная таблица экспериментальных серий}

\begin{table}[htbp]
\caption{Сводная характеристика экспериментальных серий.}
\label{tab:series_summary}
\label{tab:last}
\centering
\footnotesize
\setlength{\tabcolsep}{4pt}
\renewcommand{\arraystretch}{1.15}
\begin{tabularx}{\textwidth}{@{}c l l c >{\raggedright\arraybackslash}X@{}}
\toprule
\textbf{Серия} & \textbf{Среда} & \textbf{Метод} & \textbf{Запуск.} & \textbf{Цель} \\
\midrule
1 & LIBERO & VLA-baseline & 27 & Отправная точка SR \\
2 & LIBERO & RL-схемы (OpenVLA) & 9 & Выбор метода RL \\
3 & LIBERO & Residual SAC + stage & 27 & RL всех VLA \\
4 & LIBERO & Абляция & 5 & Вклад компонентов \\
5 & MetaWorld & VLM-код + SAC & 12 & Альтернатива без VLA \\
\bottomrule
\end{tabularx}
\end{table}

\section*{А.5. Структура репозитория прототипа}
\addcontentsline{toc}{section}{А.5. Структура репозитория прототипа}

    Код размещён в репозитории \url{https://github.com/mfclabber/bachelor-thesis}. Основные директории:
{\footnotesize
\begin{verbatim}
bachelor-thesis/
  configs/             # Hydra: series1_baseline.yaml, ...
  src/vla_rl/
    evaluate_vla.py      # инференс 9 VLA на LIBERO
    residual_rl_train.py # SAC residual-обёртка
    stage_reward.py      # 4 фазы манипуляции
    vlm_reward_gen.py    # Qwen2.5-VL (серия 5)
  scripts/
    run_series1.sh ... run_series5.sh
  results/
    libero/              # CSV SR, чекпоинты, кривые
    metaworld/           # логи, reward_code/
\end{verbatim}
}

\section*{А.6. Параметры VLM-кодогенерации (серия~5)}
\addcontentsline{toc}{section}{А.6. Параметры VLM-кодогенерации}

    Модель Qwen2.5-VL-7B-Instruct (vLLM), температура 0,7, до 2048 токенов, 1--3 итерации рефлексии. Промпт: название задачи MetaWorld, семантика наблюдений, обратная связь (SR, $\rho(R,S)$, описание поведения). Пример эволюции кода награды для push-v3:

    \textit{Итерация~1} (SR\,=\,48,3\%, $\rho = 0{,}12$):
{\footnotesize
\begin{verbatim}
def compute_dense_reward(self, action, obs):
    ee = obs["achieved_goal"][:3]
    obj = obs["observation"][:3]
    return -np.linalg.norm(ee - obj)
\end{verbatim}
}

    \textit{Итерация~3} (SR\,=\,85,0\%, $\rho = 0{,}81$):
{\footnotesize
\begin{verbatim}
def compute_dense_reward(self, action, obs):
    ee = obs["achieved_goal"][:3]
    obj = obs["observation"][:3]
    goal = obs["desired_goal"][:3]
    dist = np.linalg.norm(ee - obj)
    contact = 1.0 if abs(ee[2] - obj[2]) < 0.02 else 0.0
    d_now = np.linalg.norm(obj - goal)
    progress = self._prev_dist - d_now
    self._prev_dist = d_now
    return -dist + 0.5 * contact + 2.0 * progress
\end{verbatim}
}

\section*{А.7. Расчёт 95\% доверительного интервала}
\addcontentsline{toc}{section}{А.7. Расчёт 95\% доверительного интервала}

    Для каждой конфигурации (модель $\times$ набор LIBERO) SR усредняется по 10 задачам и 3 seed среды. 95\% ДИ вычисляется по $t$-распределению Стьюдента с $n-1 = 2$ степенями свободы:
\begin{equation*}
    \text{ДИ}_{95\%} = \bar{x} \pm t_{0{,}975,\,2} \cdot \frac{s}{\sqrt{n}}, \quad t_{0{,}975,\,2} \approx 4{,}303.
\end{equation*}
    Например, для OpenVLA после RL по трём seed: $\bar{x} = 98{,}6\%$, выборочное стандартное отклонение $s = 0{,}12\%$, $n = 3$. Тогда стандартная ошибка среднего $s/\sqrt{n} = 0{,}12/\sqrt{3} = 0{,}069\%$, а половина ширины интервала
\begin{equation*}
    t_{0{,}975,\,2}\cdot \frac{s}{\sqrt{n}} = 4{,}303 \cdot 0{,}069 = 0{,}30\%,
\end{equation*}
    откуда ДИ $= 98{,}6 \pm 0{,}3\%$. Для моделей с большим разбросом по seed (например, RT-1-X, $s \approx 0{,}32\%$) тем же способом получается более широкий интервал $\pm 0{,}8\%$.

\section*{А.8. Структура residual RL (псевдокод)}
\addcontentsline{toc}{section}{А.8. Структура residual RL}

{\footnotesize
\begin{verbatim}
# Residual RL training loop (OpenVLA + SAC)
vla = load_vla("openvla-7b-oft-libero").eval()  # frozen
residual = SAC("MlpPolicy", env, policy_kwargs=256x256)

obs, instruction = env.reset()              # reset once at start
for step in range(300_000):
    a_vla = vla.predict(obs, instruction)   # no grad
    delta = residual.predict(obs_embed)     # trainable
    action = a_vla + delta
    next_obs, reward, done = env.step(action)
    reward = stage_reward(next_obs, phase)  # 4 phases
    residual.replay_buffer.add(obs, action, reward, next_obs, done)
    obs = next_obs
    if done:                                # reset only at episode end
        obs, instruction = env.reset()
\end{verbatim}
}

\endinput
